%قبل از شروع مسئله ترجیحاً چیزی ننویسید. اما اگر خواستید بنویسید، در قالب لیستی (فقرات و شمارش) بنویسید.
%برای اضافه کردن نمره، بعدی از شروع‌مسئله، [x] را اضافه کنید و به جای x  نمره مورد نظر را قرار دهید. برای زیرمسئله‌ها هم به همین صورت است.
\شروع‌مسئله[25]
اخیراً مسعود دست از پژوهش برداشته و بازی کردن با مکعب روبیک را امری مفیدتر از کار خود یافته و فکر می‌کند حل کردن این پازل کمک بیشتری به بشریت خواهد کرد. او که اخیرا به حل مکعب روبیک اعتیاد پیدا کرده است، یک مکعب جیبیِ! 
$2 \times 2 \times 2$
خریده که وقتی بیکار است، آن را حل کند تا در مسابقه‌ای که در هفته‌های آتی دارد بتواند موفق‌تر ظاهر شود.
از این‌رو او برای حل این مساله در کمتر از ۵ ثانیه، به یاری یک کامپیوتر و صد البته کسی که بتواند کد بزند نیاز دارد، از این‌رو طبق روال، شما را به مسعود معرفی کرده‌ایم.
در کنار مستند سؤالات، یک پوشه با نام \متن‌تایپ{rubik} قرار دارد که بازی روبیک در آن پیاده‌سازی شده‌است.

هدف شما پیدا کردن کوتاه‌ترین مسیر بین حالت ابتدایی و حالت نهایی با استفاده از  الگوریتم \متن‌لاتین{BFS} است. اما استفاده از \متن‌لاتین{BFS} ساده به طوری که تا عمق ۱۴(که قطر گراف حالت‌های روبیک است) برود، کند است و مدت زمان اجرای آن به دقیقه می‌کشد و از حافظه زیادی استفاده می‌کند. حال باید الگوریتمی پیاده‌سازی کنید که مساله را در کم‌تر از ۵ ثانیه حل کند.

به‌این منظور، باید تابع \متن‌تایپ{shortest-path} که در فایل \متن‌تایپ{solver.py} است را  پیاده‌سازی کنید که دو حالت را به عنوان ورودی می‌گیرد و لیستی که شامل کوتاه‌ترین مسیر بین این دو وضعیت است را به عنوان خروجی برگرداند.

همچنین توضیحات بیشتر درباره پروژه را می‌توانید در فایل‌های \متن‌تایپ{README.txt} ای که در پوشه مربوطه هست پیدا کنید. علاوه بر این، تعدادی \متن‌لاتین{unit test} در فایل \متن‌تایپ{test-solver.py} آمده است که می‌توانید برنامه خود را با کمک آن‌ها تست کنید. نحوۀ ارسال این سوال در صفحۀ کوئرای درس قابل مشاهده است.
%اگر مسئله چند قسمت داشت به این صورت عمل کنید:
%\شروع{زیرمسائل}
%\زیرمسئله[0] % مسئلۀ ۱آ
%\زیرمسئله % مسئلۀ ۱ب
%\زیرمسئله % مسئلۀ ۱پ
%\زیرمسئله % مسئلۀ ۱ت
%\پایان{زیرمسائل}

\solution{
	فایل کد نمونه با نام \متن‌لاتین{solver.py} ضمیمه شده است.
	کد منتخب این بخش از خانم زینب احیایی انتخاب شده‌است.

}




%%%%%%%%%%%%%%%%%%%%%%%%%%%%%%%%%%%%%%%%%%%%%%%%%%%%%
% برخی دستورات مفید برای تمرین‌ها. %
%%%%%%%%%%%%%%%%%%%%%%%%%%%%%%%%%%%%%%%%%%%%%%%%%%%%%

% فرمول‌های نمونه:
%$f(x) = \Theta(x)$
%$T(x, y) \leq \log(x) + 2^y + \binom{2n}{n}$
% {\tt code\_function}


% لیست موارد بدون شماره‌گذاری:
%\شروع{فقرات}
%    \فقره مورد اول 
%        \شروع{فقرات}
%            \فقره مورد اول داخلی 
%                \شروع{فقرات}
%                    \فقره اولین مورد خیلی داخلی 
%                    \فقره دومین مورد خیلی داخلی 
%                \پایان{فقرات}
%            \فقره مورد دوم داخلی 
%        \پایان{فقرات}
%    \فقره مورد دوم 
%\پایان{فقرات}

% لیست موارد شماره‌گذاری شده:
%\شروع{شمارش}
%    \فقره یک مورد 
%    \فقره یک مورد دیگر 
%    \فقره باز هم یک مورد دیگر
%\پایان{شمارش}

% برای نوشتن یک فرمول چند خطی که یک شماره بخوره. برای حالت‌های دیگر به سایت زیر مراجعه کنید:
% https://www.overleaf.com/learn/latex/Aligning_equations_with_amsmath
%\begin{align} 
%    \begin{split}
%        (x+y)^3 &= (x+y)^2(x+y) \\
%                &= (x^2+2xy+y^2)(x+y) \\
%                &= (x^3+2x^2y+xy^2) + (x^2y+2xy^2+y^3) \\
%                &= x^3+3x^2y+3xy^2+y^3
%    \end{split}                                 
%\end{align}

% نوشتن ماتریس و آرایه دوبعدی
%\begin{align}
%    A = 
%    \begin{bmatrix}
%        A_{11} & A_{21} \\
%        A_{21} & A_{22}
%    \end{bmatrix}
%\end{align}

% اضافه کردن تصویر با شرح. برای اضافه کردن تصویر بدون شرح، فقط دستور درج‌تصویر را استفاده کنید.
%\شروع{شکل}[h!]
%\تنظیم‌ازوسط
%  \درج‌تصویر[پهنا=0.5\پهنای‌متن]{img.jpg}
%\شرح{شرح تصویر}
%\برچسب{تصویر:۱}
%\پایان{شکل}
%\پایان{زیرمسائل}

%کد زیر را ببینید.
%%\lstinputlisting{source_filename.py}
%\begin{latin}
%\begin{lstlisting}
%# Import the modules
%import sys
%import random
%
%ans = True
%
%while ans:
%question = raw_input("Ask the magic 8 ball a question: (press enter to quit) ")
%
%answers = random.randint(1,8)
%
%if question == "":
%sys.exit()
%
%elif answers == 1:
%print "It is certain"
%
%elif answers == 2:
%print "Outlook good"
%
%elif answers == 3:
%print "You may rely on it"
%
%elif answers == 4:
%print "Ask again later"
%
%elif answers == 5:
%print "Concentrate and ask again"
%
%elif answers == 6:
%print "Reply hazy, try again"
%
%elif answers == 7:
%print "My reply is no"
%
%elif answers == 8:
%print "My sources say no"
%\end{lstlisting}
%\end{latin}
